% GENERATED FILE. Edit canonical lessons, then run npm run generate:books.
% canonical-source-hash: fnv1a64:9856485b115ca3d8
% canonical-lessons: MW-C08-hear-maa, MW-C08-maa, MW-C08-hear-baap, MW-W08-ba, MW-C08-baap, MW-C08-hear-bhai, MW-C08-bhai, MW-C08-hear-bahan, MW-C08-bahan, MW-C08-family-four

\chapter{Four Family Words: Meaning First, One New Sign}
\label{ch:mw-family-four}
\hlchaptermodality{10}

\begin{chapteropening}
\textbf{By the end of this chapter:} \emph{I can recognise, say, read, and independently write four source-attested close-family labels after learning only one new Devanagari sign.}

Four close-family labels secured by ear before writing, with only \mw{ब} added and every skill scored separately.
\end{chapteropening}

\section[mā̃]{Hear \emph{mā̃} — mother}

\label{lesson:MW-C08-hear-maa}



\begin{quote}
Say the line that promises another meeting and give its two meaning parts. Then write \textbf{\mw{े}} and the known word \textbf{\mw{थारो}} once.
\end{quote}

\subsection*{What to know first}
\begin{quote}
\emph{mā̃} — \textbf{mother}
\end{quote}

Listen twice. The vowel is long and the ending is nasal. Say the whole word once; do not spell it yet. Families and the names people use for relatives vary, so this lesson teaches one source-attested everyday label, not one family model.

\subsection*{Your turn}
Point to the mother cue when you hear \emph{mā̃}, then say the word once from the meaning alone.

\begin{tcolorbox}[breakable,colback=teal!4,colframe=teal!35!black,title={Before you move on}]
Say \emph{mā̃} once, then stop.

Sources: \href{https://www.marwaripathshala.com/marwari-lesson-4-english}{Marwari Pathshala, Lesson 4} and \href{https://selindia.org/wp-content/uploads/2025/05/SWL-Marwari\_IPA.pdf}{Society for Endangered Languages, Marwari Swadesh list}.
\end{tcolorbox}
\section[mā̃]{\mw{मां} — a new word from known signs}

\label{lesson:MW-C08-maa}



\begin{quote}
Say \emph{mā̃} from its meaning. Write \textbf{\mw{म}}, \textbf{\mw{ा}}, and \textbf{\mw{ं}}; then write \textbf{\mw{ि}} and say the known question word \emph{kāĩ}.
\end{quote}

\subsection*{What to know first}
\begin{quote}
\textbf{\mw{म}} + \textbf{\mw{ा}} + \textbf{\mw{ं}} $\to$ \textbf{\mw{मां}} — \emph{mā̃} — mother
\end{quote}

Nothing on the page is new. Join three earned units and read the complete word.

\subsection*{Your turn}
Look at \textbf{\mw{मां}} for five seconds. Cover it, wait five seconds, then write the word once. Uncover and compare the long-vowel sign and nasal mark separately.

\begin{tcolorbox}[breakable,colback=teal!4,colframe=teal!35!black,title={Before you move on}]
Read and write \textbf{\mw{मां}} once without romanization.

Source: \href{https://www.marwaripathshala.com/marwari-lesson-4-english}{Marwari Pathshala, Lesson 4}.
\end{tcolorbox}
\section[bāp]{Hear \emph{bāp} — father}

\label{lesson:MW-C08-hear-baap}



\begin{quote}
Say and write mother, then say water. Write familiar \textbf{\mw{ल}} and \textbf{\mw{क}} once each.
\end{quote}

\subsection*{What to know first}
\begin{quote}
\emph{bāp} — \textbf{father}
\end{quote}

Hear one long middle vowel: \emph{bāp}. Say the word once without looking for its spelling. A respectful form such as \emph{bāpūjī} can come later; this step keeps one meaning and one sound small.

\subsection*{Your turn}
Hear \emph{mā̃} and \emph{bāp} in either order. Point to mother or father, then say only the word you heard.

\begin{tcolorbox}[breakable,colback=teal!4,colframe=teal!35!black,title={Before you move on}]
Say \emph{mā̃} and \emph{bāp} once each, then stop.

Sources: \href{https://www.marwaripathshala.com/marwari-lesson-4-english}{Marwari Pathshala, Lesson 4} and \href{https://nepal.sil.org/resources/archives/63768}{SIL International, Marwari–English Dictionary archive}.
\end{tcolorbox}
\section[ba]{\mw{ब} — one new consonant}

\label{lesson:MW-W08-ba}



\begin{quote}
Say father and mother from meaning cues. Write \textbf{\mw{पाछे} \mw{मिलसू}} once, then familiar \textbf{\mw{प}}; today's voiced sound gets a different sign.
\end{quote}

\subsection*{Script — one sign}
\begin{quote}
\textbf{\mw{ब}} — \emph{ba}
\end{quote}

Keep the model visible. Trace the lower turn, rise to the upright, and add the headline. Say \emph{ba} after the hand stops.

\subsection*{Your turn}
Copy \textbf{\mw{ब}} once. Cover it, wait five seconds, and write it once. Then say the family word \emph{bāp} that this sign will begin.

\begin{tcolorbox}[breakable,colback=teal!4,colframe=teal!35!black,title={Before you move on}]
Read \textbf{\mw{ब}} once, then stop.

Source: \href{https://www.unicode.org/charts/PDF/U0900.pdf}{Unicode Devanagari chart}.
\end{tcolorbox}
\section[bāp]{\mw{बाप} — father on the page}

\label{lesson:MW-C08-baap}



\begin{quote}
Write \textbf{\mw{ब}}, \textbf{\mw{ा}}, and \textbf{\mw{प}} separately. Say mother and father, write \textbf{\mw{कांई}}, then give the earlier see-you-later four-skill response once.
\end{quote}

\subsection*{What to know first}
\begin{quote}
\textbf{\mw{ब}} + \textbf{\mw{ा}} + \textbf{\mw{प}} $\to$ \textbf{\mw{बाप}} — \emph{bāp} — father
\end{quote}

Only the first sign is newly learned. Read the whole word without romanization.

\subsection*{Your turn}
Look for five seconds, cover \textbf{\mw{बाप}}, wait five seconds, and write it. Check that \textbf{\mw{ब}} did not turn into \textbf{\mw{प}}.

\begin{tcolorbox}[breakable,colback=teal!4,colframe=teal!35!black,title={Before you move on}]
Read \textbf{\mw{मां} — \mw{बाप}} once, then stop.

Source: \href{https://www.marwaripathshala.com/marwari-lesson-4-english}{Marwari Pathshala, Lesson 4}.
\end{tcolorbox}
\section[bhāī]{Hear \emph{bhāī} — brother}

\label{lesson:MW-C08-hear-bhai}



\begin{quote}
Say mother and write father. Ask and answer the known name exchange, then recall the formal gratitude word that begins with breathy \emph{bh}.
\end{quote}

\subsection*{What to know first}
\begin{quote}
\emph{bhāī} — \textbf{brother}
\end{quote}

Hear two long vowel parts: \emph{bhā-ī}. Say them smoothly once. The first consonant matches the familiar start of \emph{ābhār}; spelling waits for the next step.

\subsection*{Your turn}
Distinguish \emph{bāp} from \emph{bhāī} by meaning and by the breath after \emph{bh}. Answer a brother cue with \emph{bhāī} once.

\begin{tcolorbox}[breakable,colback=teal!4,colframe=teal!35!black,title={Before you move on}]
Say \emph{bhāī} once, then stop.

Source: \href{https://www.marwaripathshala.com/marwari-lesson-4-english}{Marwari Pathshala, Lesson 4}.
\end{tcolorbox}
\section[bhāī]{\mw{भाई} — brother from familiar units}

\label{lesson:MW-C08-bhai}



\begin{quote}
Write \textbf{\mw{भ}}, \textbf{\mw{ा}}, and independent \textbf{\mw{ई}}. Say mother, then ask the known polite wellbeing question with \emph{āp} and \emph{kaiso}.
\end{quote}

\subsection*{What to know first}
\begin{quote}
\textbf{\mw{भ}} + \textbf{\mw{ा}} + \textbf{\mw{ई}} $\to$ \textbf{\mw{भाई}} — \emph{bhāī} — brother
\end{quote}

The written pieces are old; only their family meaning is new.

\subsection*{Your turn}
Look at \textbf{\mw{भाई}}, cover it, wait five seconds, and write it once. Read the word without romanization and give its meaning.

\begin{tcolorbox}[breakable,colback=teal!4,colframe=teal!35!black,title={Before you move on}]
Read \textbf{\mw{भाई}} once, then stop.

Source: \href{https://www.marwaripathshala.com/marwari-lesson-4-english}{Marwari Pathshala, Lesson 4}.
\end{tcolorbox}
\section[bahan]{Hear \emph{bahan} — sister}

\label{lesson:MW-C08-hear-bahan}



\begin{quote}
Say father and brother. Write the known wellbeing question once, then \textbf{\mw{ब}}; keep the next new word oral for now.
\end{quote}

\subsection*{What to know first}
\begin{quote}
\emph{bahan} — \textbf{sister}
\end{quote}

Hear two small parts: \emph{ba-han}. Say the whole word once. A familiar alternative such as \emph{dīdī} may be used in some relationships; this lesson scores only the source form \emph{bahan}.

\subsection*{Your turn}
Hear \emph{bhāī} or \emph{bahan}. Point to brother or sister, then repeat the heard word.

\begin{tcolorbox}[breakable,colback=teal!4,colframe=teal!35!black,title={Before you move on}]
Say \emph{bhāī — bahan} once, then stop.

Source: \href{https://www.marwaripathshala.com/marwari-lesson-4-english}{Marwari Pathshala, Lesson 4}.
\end{tcolorbox}
\section[bahan]{\mw{बहन} — sister on the page}

\label{lesson:MW-C08-bahan}



\begin{quote}
Write \textbf{\mw{ब}}, \textbf{\mw{ह}}, and \textbf{\mw{न}} separately. Say brother, then give the heard wellbeing answer with \emph{hū̃} and \emph{ṭhīk}.
\end{quote}

\subsection*{What to know first}
\begin{quote}
\textbf{\mw{ब}} + \textbf{\mw{ह}} + \textbf{\mw{न}} $\to$ \textbf{\mw{बहन}} — \emph{bahan} — sister
\end{quote}

All three signs are earned. Read the word as two sound groups: \emph{ba-han}.

\subsection*{Your turn}
Hear \emph{bahan}, wait five seconds, and write \textbf{\mw{बहन}} without a visible model. Check the middle \textbf{\mw{ह}} rather than guessing from the meaning.

\begin{tcolorbox}[breakable,colback=teal!4,colframe=teal!35!black,title={Before you move on}]
Read \textbf{\mw{बहन} — \mw{भाई}} once, then stop.

Source: \href{https://www.marwaripathshala.com/marwari-lesson-4-english}{Marwari Pathshala, Lesson 4}.
\end{tcolorbox}
\section[Practice]{Checkpoint — four close-family words}

\label{lesson:MW-C08-family-four}



\begin{quote}
Write \textbf{\mw{ू}} once and say its long vowel. Then close the earlier page.
\end{quote}

\subsection*{What to know first}
\begin{quote}
\textbf{\mw{मां} — \mw{बाप} — \mw{भाई} — \mw{बहन}}
\end{quote}

The four labels name relationships; they do not prescribe one family structure.

\subsection*{Your turn}
1. \textbf{Listen:} identify all four heard labels. 2. \textbf{Speak:} produce two different labels from meanings. 3. \textbf{Read:} match all four printed words. 4. \textbf{Write:} write two heard labels without a model.

Score each skill separately.

\begin{tcolorbox}[breakable,colback=teal!4,colframe=teal!35!black,title={Before you move on}]
Stop after one independently accurate pass.

Source: \href{https://www.marwaripathshala.com/marwari-lesson-4-english}{Marwari Pathshala, Lesson 4}.
\end{tcolorbox}
